% Balanced complete-edition layout; entries unchanged from the accepted bibliography.
\begin{thebibliography}{99}
\label{refs:chapter07}
\bibitem[Black and Wiliam(1998)]{black1998}Black, P., and Wiliam, D. (1998).
Assessment and classroom learning.
\emph{Assessment in Education: Principles, Policy \& Practice}, 5(1), 7--74.
\doi{10.1080/0969595980050102}.

\bibitem[Buffardi(2020)]{buffardi2020}Buffardi, K. (2020).
Assessing individual contributions to software engineering projects with Git logs and user stories.
In \emph{Proceedings of the 51st ACM Technical Symposium on Computer Science Education}, pp.~650--656.
\doi{10.1145/3328778.3366948}.

\bibitem[Course study manuscript(n.d.)]{seetmanuscript}Course study manuscript (n.d.).
\emph{Assessing Student Contributions with Network-Based Evidence in an Open Source Software Engineering Course}.
Unpublished anonymous manuscript, 11 pp. Author-supplied version received 1 November 2026; reports a Spring 2025 deployment.

\bibitem[Espeland and Sauder(2007)]{espeland2007}Espeland, W.~N., and Sauder, M. (2007).
Rankings and reactivity: How public measures recreate social worlds.
\emph{American Journal of Sociology}, 113(1), 1--40.
\doi{10.1086/517897}.

\bibitem[Forsgren et~al.(2021)]{forsgren2021}Forsgren, N., Storey, M.-A., Maddila, C., Zimmermann, T., Houck, B., and Butler, J. (2021).
The SPACE of developer productivity: There's more to it than you think.
\emph{ACM Queue}, 19(1), 20--48.
\doi{10.1145/3454122.3454124}.

\bibitem[Han et~al.(2026a)]{han2026rank}Han, F., Zhao, S., Wang, W., Peng, J., and Xia, X. (2026a).
OpenRank: A centrality algorithm for high-dimensional heterogeneous networks in open source collaboration.
\emph{Information Sciences}, 730, 122893.
\doi{10.1016/j.ins.2025.122893}.

\bibitem[Hundhausen et~al.(2022)]{hundhausen2022}Hundhausen, C.~D., Conrad, P.~T., Carter, A.~S., and Adesope, O. (2022).
Assessing individual contributions to software engineering projects: A replication study.
\emph{Computer Science Education}, 32(3), 335--354.
\doi{10.1080/08993408.2022.2071543}.

\bibitem[Hundhausen et~al.(2023)]{hundhausen2023}Hundhausen, C., Conrad, P., Adesope, O., and Tariq, A. (2023).
Combining GitHub, chat, and peer evaluation data to assess individual contributions to team software development projects.
\emph{ACM Transactions on Computing Education}, 23(3), Article 33, 23 pp.
\doi{10.1145/3593592}.

\bibitem[Kalliamvakou et~al.(2014)]{kalliamvakou2014}Kalliamvakou, E., Gousios, G., Blincoe, K., Singer, L., German, D.~M., and Damian, D. (2014).
The promises and perils of mining GitHub.
In \emph{Proceedings of the 11th Working Conference on Mining Software Repositories}, pp.~92--101.
\doi{10.1145/2597073.2597074}.

\bibitem[Kane(2013)]{kane2013}Kane, M.~T. (2013).
Validating the interpretations and uses of test scores.
\emph{Journal of Educational Measurement}, 50(1), 1--73.
\doi{10.1111/jedm.12000}.

\bibitem[Kluger and DeNisi(1996)]{kluger1996}Kluger, A.~N., and DeNisi, A. (1996).
The effects of feedback interventions on performance: A historical review, a meta-analysis, and a preliminary feedback intervention theory.
\emph{Psychological Bulletin}, 119(2), 254--284.
\doi{10.1037/0033-2909.119.2.254}.

\bibitem[Maruping et~al.(2019)]{maruping2019}Maruping, L.~M., Daniel, S.~L., and Cataldo, M. (2019).
Developer centrality and the impact of value congruence and incongruence on commitment and code contribution activity in open source software communities.
\emph{MIS Quarterly}, 43(3), 951--976.
\url{https://aisel.aisnet.org/misq/vol43/iss3/14/}.

\end{thebibliography}
\clearpage
\begingroup
\renewcommand{\bibsection}{\section*{References\enspace{\normalfont\sffamily\fontsize{11}{14}\selectfont continued}}}
\begin{thebibliography}{99}
\bibitem[Messick(1994)]{messick1994}Messick, S. (1994).
Validity of psychological assessment: Validation of inferences from persons' responses and performances as scientific inquiry into score meaning.
\emph{ETS Research Report Series}, 1994(2), i--28.
\doi{10.1002/j.2333-8504.1994.tb01618.x}.

\bibitem[Nicol and Macfarlane-Dick(2006)]{nicol2006}Nicol, D.~J., and Macfarlane-Dick, D. (2006).
Formative assessment and self-regulated learning: A model and seven principles of good feedback practice.
\emph{Studies in Higher Education}, 31(2), 199--218.
\doi{10.1080/03075070600572090}.

\bibitem[Pinto et~al.(2019)]{pinto2019}Pinto, G., Ferreira, C., Souza, C., Steinmacher, I., and Meirelles, P. (2019).
Training software engineers using open-source software: The students' perspective.
In \emph{Proceedings of the IEEE/ACM 41st International Conference on Software Engineering: Software Engineering Education and Training}, pp.~147--157.
\url{https://www.igor.pro.br/publica/papers/ICSE_SEET2019.pdf}.

\bibitem[Sadler(1989)]{sadler1989}Sadler, D.~R. (1989).
Formative assessment and the design of instructional systems.
\emph{Instructional Science}, 18, 119--144.
\doi{10.1007/BF00117714}.

\bibitem[Steinmacher et~al.(2015)]{steinmacher2015}Steinmacher, I., Conte, T.~U., Gerosa, M.~A., and Redmiles, D.~F. (2015).
Social barriers faced by newcomers placing their first contribution in open source software projects.
In \emph{Proceedings of the 18th ACM Conference on Computer Supported Cooperative Work \& Social Computing}, pp.~1379--1392.
\doi{10.1145/2675133.2675215}.

\bibitem[Wang et~al.(2024)]{wang2024course}Wang, J., Huang, W., Zhao, S., Xia, X., Han, F., Wang, W., and Zhang, Y. (2024).
OpenRank contribution evaluation method and empirical study in open-source course.
\emph{Journal of East China Normal University (Natural Science)}, 2024(5), 11--19. In Chinese.
\doi{10.3969/j.issn.1000-5641.2024.05.002}.

\bibitem[X-lab(n.d.)]{xlabeducation}X-lab (n.d.).
\emph{Making Contributions Visible and Providing Pathways for Growth: East China Normal University's Integrated Open Source Development Platform}.
Author-supplied institutional practice account, in Chinese, received 1 November 2026. English title translated for this book; not an independently evaluated outcome study.

\bibitem[Zagalsky et~al.(2015)]{zagalsky2015}Zagalsky, A., Feliciano, J., Storey, M.-A., Zhao, Y., and Wang, W. (2015).
The emergence of GitHub as a collaborative platform for education.
In \emph{Proceedings of the 18th ACM Conference on Computer Supported Cooperative Work \& Social Computing}, pp.~1906--1917.
\url{https://alexeyza.com/pdf/cscw15.pdf}.

\bibitem[Zhao et~al.(2024)]{zhao2024}Zhao, S., Xia, X., Fitzgerald, B., Li, X., Lenarduzzi, V., Taibi, D., Wang, R., Wang, W., and Tian, C. (2024).
OpenRank Leaderboard: Motivating open source collaborations through social network evaluation in Alibaba.
In \emph{Proceedings of the 46th International Conference on Software Engineering: Software Engineering in Practice}, pp.~346--357.
\doi{10.1145/3639477.3639734}.
\end{thebibliography}
\endgroup
